% GENERATED FILE. Edit canonical lessons, then run npm run generate:books.
% canonical-source-hash: fnv1a64:2bf6a4b3a5037d49
% canonical-lessons: MR-C239-nikal, MR-C239-pharak, MR-C239-udaharan, MR-C239-satya, MR-C239-niyam

\chapter{Result, Difference, Example, Truth, Rule}
\label{ch:mr-result-difference-example-truth-rule}
\hlchaptermodality{239}

\begin{chapteropening}
\textbf{By the end of this chapter:} \emph{I can say a result, a difference, an example, the truth and a rule.}

Complete the last lesson of chapter 239: I can say a result, a difference, an example, the truth and a rule.
\end{chapteropening}

\section[\texorpdfstring{\mr{निकाल} (nikāl)}{निकाल (nikāl)}]{\mr{निकाल} (nikāl) — a result}

\label{lesson:MR-C239-nikal}



\begin{quote}
Before the new one: say the Marathi for a tourist, then the Marathi for an adult.
\end{quote}

\subsection*{You'll want to know: \mr{निकाल}}
\textbf{\mr{निकाल}} — \emph{nikāl} — \textquotedblleft{}a result\textquotedblright{}.

\begin{grammarlens}[title={What you've built}]
One more word for this chapter.
\end{grammarlens}

\subsection*{Your turn}
\begin{itemize}
\raggedright
  \item \emph{Say it:} \emph{nikāl}
  \item \emph{Say it:} \emph{nikāl}, once more
  \item \emph{Say it:} say \emph{nikāl}
  \item \emph{Recall:} say the Marathi for a tourist, then the Marathi for an adult, then say \emph{nikāl} again
\end{itemize}

\begin{tcolorbox}[breakable,colback=teal!4,colframe=teal!35!black,title={Before you move on}]
What does \mr{निकाल} mean? (\textquotedblleft{}A result\textquotedblright{}.) Say it once more.
\end{tcolorbox}
\section[\texorpdfstring{\mr{फरक} (pharak)}{फरक (pharak)}]{\mr{फरक} (pharak) — a difference}

\label{lesson:MR-C239-pharak}



\begin{quote}
Before the new one: say the Marathi for an adult, then the Marathi for a result.
\end{quote}

\subsection*{You'll want to know: \mr{फरक}}
\textbf{\mr{फरक}} — \emph{pharak} — \textquotedblleft{}a difference\textquotedblright{}.

\begin{grammarlens}[title={What you've built}]
One more word for this chapter.
\end{grammarlens}

\subsection*{Your turn}
\begin{itemize}
\raggedright
  \item \emph{Say it:} \emph{pharak}
  \item \emph{Say it:} \emph{pharak}, once more
  \item \emph{Say it:} say \emph{pharak}
  \item \emph{Recall:} say the Marathi for an adult, then the Marathi for a result, then say \emph{pharak} again
\end{itemize}

\begin{tcolorbox}[breakable,colback=teal!4,colframe=teal!35!black,title={Before you move on}]
What does \mr{फरक} mean? (\textquotedblleft{}A difference\textquotedblright{}.) Say it once more.
\end{tcolorbox}
\section[\texorpdfstring{\mr{उदाहरण} (udāharaṇ)}{उदाहरण (udāharaṇ)}]{\mr{उदाहरण} (udāharaṇ) — an example}

\label{lesson:MR-C239-udaharan}



\begin{quote}
Before the new one: say the Marathi for a result, then the Marathi for a difference.
\end{quote}

\subsection*{You'll want to know: \mr{उदाहरण}}
\textbf{\mr{उदाहरण}} — \emph{udāharaṇ} — \textquotedblleft{}an example\textquotedblright{}.

\begin{grammarlens}[title={What you've built}]
One more word for this chapter.
\end{grammarlens}

\subsection*{Your turn}
\begin{itemize}
\raggedright
  \item \emph{Say it:} \emph{udāharaṇ}
  \item \emph{Say it:} \emph{udāharaṇ}, once more
  \item \emph{Say it:} say \emph{udāharaṇ}
  \item \emph{Recall:} say the Marathi for a result, then the Marathi for a difference, then say \emph{udāharaṇ} again
\end{itemize}

\begin{tcolorbox}[breakable,colback=teal!4,colframe=teal!35!black,title={Before you move on}]
What does \mr{उदाहरण} mean? (\textquotedblleft{}An example\textquotedblright{}.) Say it once more.
\end{tcolorbox}
\section[\texorpdfstring{\mr{सत्य} (satya)}{सत्य (satya)}]{\mr{सत्य} (satya) — the truth}

\label{lesson:MR-C239-satya}



\begin{quote}
Before the new one: say the Marathi for a difference, then the Marathi for an example.
\end{quote}

\subsection*{You'll want to know: \mr{सत्य}}
\textbf{\mr{सत्य}} — \emph{satya} — \textquotedblleft{}the truth\textquotedblright{}.

\begin{grammarlens}[title={What you've built}]
One more word for this chapter.
\end{grammarlens}

\subsection*{Your turn}
\begin{itemize}
\raggedright
  \item \emph{Say it:} \emph{satya}
  \item \emph{Say it:} \emph{satya}, once more
  \item \emph{Say it:} say \emph{satya}
  \item \emph{Recall:} say the Marathi for a difference, then the Marathi for an example, then say \emph{satya} again
\end{itemize}

\begin{tcolorbox}[breakable,colback=teal!4,colframe=teal!35!black,title={Before you move on}]
What does \mr{सत्य} mean? (\textquotedblleft{}The truth\textquotedblright{}.) Say it once more.
\end{tcolorbox}
\section[\texorpdfstring{\mr{नियम} (niyam)}{नियम (niyam)}]{\mr{नियम} (niyam) — a rule}

\label{lesson:MR-C239-niyam}



\begin{quote}
Before the new one: say the Marathi for an example, then the Marathi for the truth.
\end{quote}

\subsection*{You'll want to know: \mr{नियम}}
\textbf{\mr{नियम}} — \emph{niyam} — \textquotedblleft{}a rule\textquotedblright{}.

\begin{grammarlens}[title={What you've built}]
One more word for this chapter.
\end{grammarlens}

\subsection*{Your turn}
\begin{itemize}
\raggedright
  \item \emph{Say it:} \emph{niyam}
  \item \emph{Say it:} \emph{niyam}, once more
  \item \emph{Say it:} say \emph{niyam}
  \item \emph{Recall:} say the Marathi for an example, then the Marathi for the truth, then say \emph{niyam} again
\end{itemize}

\begin{tcolorbox}[breakable,colback=teal!4,colframe=teal!35!black,title={Before you move on}]
What does \mr{नियम} mean? (\textquotedblleft{}A rule\textquotedblright{}.) Say it once more.
\end{tcolorbox}
